% !TEX root = ../../main.tex
%% Corrigé du défi du chapitre 2 — vérifié par verifications/defis_01_05.py.

\begin{blocLivre}{cor-2-25a}
\begin{corrige}{25}{Sommes de puissances et théorème de Nicomaque}
\demarche{Un télescopage développé par le binôme fournit une relation de
récurrence entre les $S_p$ ; les figures en donnent la lecture
géométrique.}
\partieExo{A}{La méthode de Pascal}
\begin{questions}
  \item La somme télescope : elle vaut $(n+1)^{p+1}-1$. D'autre part,
    par la formule du binôme,
    $(k+1)^{p+1}-k^{p+1}=\sum_{j=0}^{p}\binom{p+1}{j}k^{j}$, et en
    sommant sur $k$ :
    \[
    \res{\sum_{j=0}^{p}\binom{p+1}{j}S_j(n)=(n+1)^{p+1}-1} .
    \]
  \item $p=1$ : $n+2T_n=(n+1)^{2}-1=n^{2}+2n$, d'où
    $T_n=\dfrac{n(n+1)}{2}$. $p=2$ :
    $n+3T_n+3S_2(n)=n^{3}+3n^{2}+3n$, d'où
    $3S_2(n)=n^{3}+\tfrac32n^{2}+\tfrac12n=\tfrac{n(n+1)(2n+1)}{2}$ et
    $\res{S_2(n)=\dfrac{n(n+1)(2n+1)}{6}}$.
  \item $p=3$ : $4S_3(n)=(n+1)^{4}-(n+1)-4T_n-6S_2(n)$. Or
    $(n+1)^{4}-(n+1)=n(n+1)(n^{2}+3n+3)$ et
    $4T_n+6S_2(n)=n(n+1)(2n+3)$, d'où
    $4S_3(n)=n(n+1)(n^{2}+n)=n^{2}(n+1)^{2}$ :
    $\res{S_3(n)=\Bigl(\dfrac{n(n+1)}{2}\Bigr)^{2}}$.
\end{questions}
\partieExo{B}{Le théorème de Nicomaque, en images}
\begin{questions}[resume]
  \item $T_k-T_{k-1}=k$ et $T_k+T_{k-1}=\dfrac{k(k+1)+(k-1)k}{2}=k^{2}$,
    donc $T_k^{2}-T_{k-1}^{2}=k\cdot k^{2}=k^{3}$. Géométriquement, la
    $k$-ième équerre est la différence des carrés de côtés $T_k$ et
    $T_{k-1}$ : sa largeur est $k$, son aire $k^{3}$. En sommant
    (télescopage), $\res{S_3(n)=T_n^{2}}$.
  \item Les $k$ impairs annoncés sont symétriques autour de leur moyenne
    $\dfrac{(k^{2}-k+1)+(k^{2}+k-1)}{2}=k^{2}$ : leur somme vaut
    $k\cdot k^{2}=k^{3}$. Les blocs se succèdent sans trou ($k(k+1)-1$ est
    suivi de $(k+1)k+1$) : $1^{3}+\dots+n^{3}$ est la somme des
    $1+2+\dots+n=T_n$ premiers impairs, soit $T_n^{2}$.
\end{questions}
\end{corrige}
\end{blocLivre}

\begin{blocLivre}{cor-2-25b}
\begin{corrigeSuite}
\partieExo{C}{Le cas général}
\begin{questions}[start=6]
  \item Pour $p=0$, $S_0(n)=n$ convient. Supposons le résultat vrai pour
    tous les $j<p$. La relation de la question 1 s'écrit
    \[
    (p+1)S_p(n)=\bigl((n+1)^{p+1}-1\bigr)
    -\sum_{j=0}^{p-1}\binom{p+1}{j}S_j(n) .
    \]
    Le premier terme est un polynôme en $n$ de degré $p+1$, de coefficient
    dominant $1$ et de terme constant nul ; chaque $S_j$ ($j<p$) est un
    polynôme de degré $j+1\leqslant p$ sans terme constant. Donc
    $(p+1)S_p(n)$ est polynomial de degré $p+1$, de coefficient dominant
    $1$, sans terme constant : $S_p(n)$ a pour coefficient dominant
    $\res{\dfrac{1}{p+1}}$.
  \item $p=4$ :
    $5S_4(n)=(n+1)^{5}-1-n-5T_n-10S_2(n)-10S_3(n)$ ; le calcul, puis la
    factorisation, donnent
    \[
    \res{S_4(n)=\frac{n(n+1)(2n+1)(3n^{2}+3n-1)}{30}} .
    \]
    Contrôles : $n=1$ donne $\frac{1\cdot2\cdot3\cdot5}{30}=1$ ; $n=2$
    donne $\frac{2\cdot3\cdot5\cdot17}{30}=17=1+16$.
  \item D'après la question 6,
    $S_p(n)=\dfrac{n^{p+1}}{p+1}+(\text{termes de degré}\leqslant p)$,
    donc
    \[
    \frac{1}{n}\sum_{k=1}^{n}\Bigl(\frac kn\Bigr)^{p}
    =\frac{S_p(n)}{n^{p+1}}\xrightarrow[n\to+\infty]{}\res{\frac{1}{p+1}}
    =\int_{0}^{1}t^{p}\,\mathrm{d}t .
    \]
    L'aire des rectangles de hauteurs $\bigl(\frac kn\bigr)^{p}$ tend vers
    l'aire sous la courbe de $t\mapsto t^{p}$ : c'est la méthode des
    rectangles, que le chapitre 15 généralisera à toute fonction
    continue.
\end{questions}
\end{corrigeSuite}
\end{blocLivre}
